% 核心结论：准备例：有效值相量与峰值向量
% 图原型：纯几何／对比；按本题相位与电路拓扑绘制。
\documentclass[tikz,border=7pt]{standalone}
\usepackage{fontspec,xeCJK,amsmath,amssymb}
\IfFontExistsTF{Times New Roman}{\setmainfont{Times New Roman}}{\setmainfont{TeX Gyre Termes}}
\IfFontExistsTF{Heiti SC}{\setCJKmainfont[AutoFakeBold=1.5]{Heiti SC}}{\setCJKmainfont[AutoFakeBold=1.5]{FandolHei-Regular}}
\usetikzlibrary{arrows.meta,calc,angles,quotes}
\definecolor{ink}{HTML}{243745}
\definecolor{blue}{HTML}{275DAB}
\definecolor{green}{HTML}{177C63}
\definecolor{red}{HTML}{BD4935}
\definecolor{purple}{HTML}{715393}
\begin{document}
\begin{tikzpicture}[font=\fontsize{9.5}{12}\selectfont,text=ink,
 v/.style={line width=1.1pt,-{Stealth[length=2.4mm,width=1.5mm]}},
 axis/.style={ink!65,line width=.45pt,-{Stealth[length=1.7mm,width=1.1mm]}},
 aux/.style={ink!40,densely dashed,line width=.5pt},
 lab/.style={fill=white,inner sep=2pt},
 card/.style={draw=blue!25,fill=blue!3,rounded corners=3pt,text width=6cm,align=left,inner sep=10pt}]
\path[use as bounding box] (-8,-3.35) rectangle (8,3.35);
\node[font=\fontsize{11}{14}\selectfont\bfseries] at (0.00000,3.04000) {准备例：有效值相量与峰值向量};
\begin{scope}[shift={(-4,0)}]
\draw[axis] (-2.50000,0.00000) -- (2.60000,0.00000);
\draw[axis] (0.00000,-2.05000) -- (0.00000,2.25000);
\node[] at (2.58000,-0.26000) {$\operatorname{Re}$};
\node[] at (-0.33000,2.20000) {$\operatorname{Im}$};
\node[] at (-0.18000,-0.20000) {$O$};
\draw[v,red] (0.00000,0.00000) -- (2.44949,1.41421);
\draw[v,blue] (0.00000,0.00000) -- (1.73205,1.00000);
\node[lab,text=blue] at (1.65000,0.65000) {$P$};
\node[lab,text=red] at (2.54949,1.81421) {$Q\;(U_m=10\sqrt2)$};
\draw[green] (.65,0) arc[start angle=0,end angle=30,radius=.65];
\node[lab,text=green] at (1.30000,0.28000) {$30^\circ$};
\draw[aux] (2.44949,1.41421) -- (2.44949,0.00000);
\end{scope}
\begin{scope}[shift={(4,0)}]
\node[card] at (0.00000,1.10000) {$\widetilde U=10\angle(30^\circ)\,\mathrm V$\\[4pt]$OP$：静止有效值相量。};
\node[card] at (0.00000,-1.00000) {$OQ=\sqrt2\,OP$\\[4pt]$t=0$ 时，峰值向量投影为 $u(0)$。};
\end{scope}
\node[text width=15cm,align=center] at (0.00000,-3.00000) {蓝色长度表示有效值；红色外端对应峰值。角度均从正实轴起算。};
\end{tikzpicture}
\end{document}
